% !TeX root = ../教材题库_全部条目_转写工作稿.tex
% Source: XJ-B1-C06-EX-P236-Y1549
% 第6章 / 6.4 用样本估计总体 / 教材 p.236、237 / PDF 第243、244页
% OCR未校对：公式与文字均可能有误；未完成数学TeX转写；禁止直接发给学生
% 原题及相邻原解截图（核对用）：../题目原图/XJ-B1-C06-EX-P236-Y1549-01.webp, ../题目原图/XJ-B1-C06-EX-P236-Y1549-02.webp
% 本文件只在转写完成后作为题干使用；OCR 条目尚不是可靠题干。
\begin{bankitem}{6.4 例题 例题2}{第6章 / 6.4 用样本估计总体 / 教材 p.236、237 / PDF 第243、244页}{XJ-B1-C06-EX-P236-Y1549}{OCR未校对：公式与文字均可能有误；未完成数学TeX转写；禁止直接发给学生}
\begin{ocrdraft}
[教材p.236 / PDF第243页]
\par
00表6一4是某地统计局调查100户家庭月均用水量（单位：t）的频率分布
\par
表，试根据该表估计该地每户家庭的月均用水量．
\par
[教材p.237 / PDF第244页]
\par
表6一4
\par
分组
\par
频
\par
率
\par
0．5）
\par
4
\par
0．04
\par
8
\par
0．08
\par
0，
\par
1.5）
\par
0．15
\par
[1.5，2）
\par
22
\par
0．22
\par
[2，
\par
2．5）
\par
0．25
\par
[2．5，3）
\par
14
\par
0．14
\par
[3，3．5）
\par
6
\par
0．06
\par
[3．5，4）
\par
4
\par
0．04
\par
[4，
\par
4．5）
\par
2
\par
0．02
\par
100
\par
1.00
\par
分析要确定这100户家庭的月均用水量，就必须计算其总用水量．由于每组
\par
中每户月用水量只是一个范围，因此可用各组区间的组中值（位于各组中央的值）近
\par
似地表示．
\par
（方法一）100户家庭的月总用水量约为
\par
解
\par
0．25\ensuremath{\times}4十0．75\ensuremath{\times}8十1.25\ensuremath{\times}15十1.75\ensuremath{\times}22十2．25\ensuremath{\times}25十2．75\ensuremath{\times}14十3．25\ensuremath{\times}6十
\par
3．75\ensuremath{\times}4十4．25\ensuremath{\times}2=202（t），
\par
又202\ensuremath{\div}100=2．02（t），
\par
因此，估计该地每户家庭的月均用水量为2．02t.
\par
（方法二）求组中值与对应频率之积的和．
\par
例2在计算平均数
\par
=0．25\ensuremath{\times}0．04十0．75\ensuremath{\times}0．08十1.25\ensuremath{\times}0．15十1.75\ensuremath{\times}
\par
时，用各组的组中值代表
\par
各组的实际数据，使用组
\par
0．22十2．25\ensuremath{\times}0．25十2．75\ensuremath{\times}0．14十3．25\ensuremath{\times}0．06十
\par
中值进行计算的前提是：
\par
3．75\ensuremath{\times}0，04十4．25\ensuremath{\times}0．02
\par
各组数据在组内的分布是
\par
=2．02（t），
\par
均匀的．
\par
因此，估计该地每户家庭的月均用水量为2．02t.
\par
一般地，若取值为
\par
\ensuremath{\cdot}\ensuremath{\cdot}，的频率分别为，了2，
\par
，，则其平均数
\par
为十丆2十，\ensuremath{\cdot}十丆，，
\end{ocrdraft}
\end{bankitem}
